
\usepackage{amstext}
\usepackage{stmaryrd}
%\usepackage{listings}
%\lstset{language=C++}
%\usepackage[boxed]{algorithm2e}
\usepackage{algpseudocode}
\usepackage[algo2e,linesnumbered,vlined,commentsnumbered,ruled]{algorithm2e}

\usepackage[normalem]{ulem}
\usepackage{graphicx}

%\usepackage{authblk}
%\usepackage{xcolor}
\usepackage{array}
\usepackage{verbatim}
\usepackage{booktabs}
\usepackage{calc}
\usepackage{textcomp}
\usepackage{multirow}
%\usepackage[boxed]{algorithm2e}

\usepackage[all]{xy}

\usepackage{cancel}
\usepackage{mathtools}
\usepackage{placeins}
\usepackage{xspace}

%\input{packages}

\usepackage{tikz,tikzscale}
\usetikzlibrary{calc}
\usetikzlibrary{shadings}
%\usepackage[backend=biber,style=alphabetic]{biblatex}
%\addbibresource{literature.bib}

%%%%%%%%%% Packages to help editing %%%%%%%%%%
%\usepackage[notref,notcite]{showkeys} % Show labels.
\usepackage[mode=multiuser,status=draft]{fixme} % Provide note, warning, error, fatal.
%\fxsetup{layout=inline}%pdfcsignote}
\fxsetup{envlayout=color}
\fxsetup{innerlayout=inline}
\fxsetup{targetlayout=colorcb}
\FXProvidesTargetLayout{color}

\fxusetheme{color}
\FXRegisterAuthor{ss}{envss}{SS}
\FXRegisterAuthor{mr}{envmr}{MR}
\FXRegisterAuthor{mw}{envmw}{MW}
\FXRegisterAuthor{mk}{envmk}{MK}

\usepackage{subcaption}
\usepackage[format=plain,indention=.5cm]{caption}

%%%% mathematica code
\usepackage{tcolorbox}
\usepackage{listings}
\tcbuselibrary{listings}

\definecolor{mathematicablue}{RGB}{0,0,150}
\lstdefinestyle{mathematica}{
	language=Mathematica,
	basicstyle=\ttfamily\small,
	keywordstyle=\color{mathematicablue}\bfseries,
	commentstyle=\color{gray}\itshape,
	stringstyle=\color{purple},
	showstringspaces=false,
	breaklines=true
}

%%%%%

\begin{comment}

\DontPrintSemicolon
\SetKwComment{Comment}{$\triangleleft$ }{}
\SetKwProg{Fn}{Function}{}{end}

\providecommand{\keywords}[1]
{
	\small
	\textbf{\textit{Keywords:}} #1
}
\end{comment}
% \newcommand{\FigureSubst}[1]{ \fbox{ 
% \begin{minipage} {0.7\textwidth} #1  
% \end{minipage}  }   }

\def\mesh{\mathcal{T}}
\def\cell{K}
\def\R{\mathbb R}
\def\ncol{N_{\text{colors}}}
\def\rest#1{I_{#1}^\downarrow}
\def\prol#1{I_{#1}^\uparrow}

\newcommand{\pluseq}{\stackrel{+}{\gets}}
%\newcommand{\pluseq}{\stackbin{+}{=}}

\newcommand{\Dd}{\text{d}}
\newcommand{\DD}{\text{D}}

% Mohsen's setting and definitions
%###################################%
\usepackage{bm,upgreek}
\newcommand\rd{{\rm{d}}}
\newcommand{\bfm}[1]{{\rm\bf #1}}

\newcommand{\greentext}[1]{{\color{darkgreen} #1}}
\newcommand{\bluetext}[1]{{\color{blue} #1}}
\newcommand{\redtext}[1]{{\color{red} #1}}

\DeclareMathOperator{\tr}{tr}
%###################################%

% MW settings
%###################################%

%\newcommand{\dlt}{\bm{\delta_{u}}}
%\newcommand{\corr}{\bm{{v}}}

\newcommand{\fem}[1]{\mbox{\bf \texttt{#1}}}

\newcommand{\dlt}{\, \delta {\bfm{u}}}
\newcommand{\corr}{  \,\Delta \bfm{u}\, }
% \newcommand{\resultFE}{  \, \bfm{w}\, }
\newcommand{\testFunc}{\bm{\phi}_i}
\newcommand{\inputFE}{\bfm{v}}
\newcommand{\currFE}{\, \bar{ \bfm{u}}\, }

%\newcommand{\currReal}{\, \bar{ \mathsf{u}}\, }
%\newcommand{\corrReal}{  \,\Delta \mathsf{u}\, }
%\newcommand{\resultReal}{  \, \mathsf{w}\, }
%\newcommand{\inputReal}{\mathsf{v}}
\newcommand{\currReal}{\bar{ \fem{U}}}
\newcommand{\corrReal}{\Delta \fem{U}}
\newcommand{\resultReal}{\fem{W}}
\newcommand{\inputReal}{\fem{V}}

\newcommand{\scalar }{\texttt{scalar}\xspace}
\newcommand{\scalarRef}{\texttt{scalar\_ref}\xspace}
\newcommand{\tensorS}{\texttt{tensor4}\xspace}
\newcommand{\pade}{\texttt{Pade log}\xspace}

\newcommand{\none}{\texttt{ADrecompute}\xspace}
\newcommand{\tensorAD}{\texttt{ADstore}\xspace}
\newcommand{\tensorFAST}{\texttt{tensor2}\xspace}
\newcommand{\scalarCurr}{\texttt{scalar}\xspace}

\newcommand{\crap}{\texttt{naive\_AD}\xspace}

\newcommand{\noneSplit}{\texttt{ADrecompute}\xspace}
\newcommand{\tensorADsplit}{\texttt{ADstore}\xspace}

\newcommand{\gz}{\mathbf{g}}

%\newcommand{\Grad}{\nabla}
%\newcommand{\grad}{\nabla}
\DeclareMathOperator{\Grad}{Grad}
\DeclareMathOperator{\grad}{grad}
\newcommand{\grads}{\grad^{\rm s}\!}
%###################################%
%

\usepackage{graphicx} % Required for inserting images

% \newcommand{\revOld}[1]{}
% \newcommand{\revNew}[1]{{{#1}}}
\newcommand{\revNew}[1]{{\color{blue}{#1}}}   % only in original submission
\newcommand{\revOld}[1]{\sout{\protect\color{red}{#1}}}  % only in rev 1

%% TEMPORARY - until we learn how to make lower-case \mathbb...
% \newcommand{\bbc}{{\rm c}\hspace{-0.25em}{\rm c}}
%lowercase blackboard characters :)
\usepackage{bbm}
\newcommand{\bbc}{\mathbbm{c}}

\newcommand{\cNH}{\textit{cNH}\xspace}
\newcommand{\sNH}{\textit{sNH}\xspace}

\DeclareRobustCommand{\mmodels}{\mathrel{||}\joinrel \Relbar}

\usepackage{amsmath} % Good practice for math
\usepackage{amssymb} % Provides \models
\usepackage{bm}      % Provides \bm for better bold math
\usepackage{graphicx} % Provides \scalebox

% --- Optional: Define a custom command ---
% Choose one definition based on your preferred method/size
% Method 1: Using scalebox (adjust 0.8 to your desired size factor)
\newcommand{\smbldmodels}{\mathrel{\scalebox{0.8}{\bm{\models}}}}

\usepackage{graphicx} % Required for inserting images
% \usepackage[hidelinks]{hyperref}

% simple change-tracking commands for revision rounds
\usepackage{xcolor} % safe lightweight color support

\definecolor{rev@add}{RGB}{0,0,250}   % blue for additions
\definecolor{rev@del}{RGB}{150,0,0}   % red for deletions
\definecolor{rev@note}{RGB}{120,0,120}% purple for notes

\newif\ifshowchanges
	% \showchangestrue
	\showchangesfalse  % to hide markup and produce "clean" text

	% \add{new text}          -> shows new text in blue (or plain text when hidden)
	% \del{old text}          -> shows old text struck out in red (or nothing when hidden)
	% \chg{old}{new}         -> shows old struck out in red followed by new in blue (or only new when hidden)
	% \revnote{note text}     -> small inline purple note (hidden when markup off)
	\newcommand{\add}[1]{\ifshowchanges\textcolor{rev@add}{#1}\else#1\fi}
	\newcommand{\del}[1]{\ifshowchanges\textcolor{rev@del}{\sout{#1}}\else\relax\fi}
	\newcommand{\chg}[2]{\ifshowchanges\textcolor{rev@del}{\sout{#1}}\ \textcolor{rev@add}{#2}\else#2\fi}
	\newcommand{\revnote}[1]{\ifshowchanges{\footnotesize\textcolor{rev@note}{[#1]}}\fi}

% quick toggles
% use \showchangesfalse to compile final clean version
% use \showchangestrue to show markup

\newcommand{\modelsSym}{\mathrel{\scriptstyle\bm{\models}}}
% Turnstile symbol (⊢) as a proper math relation
\newcommand{\turnstile}{\mathrel{\vdash}}

% Optional aliases / helper
\newcommand{\proves}{\turnstile}
\newcommand{\judgment}[2]{#1\ \turnstile\ #2}
